\documentclass[t,aspectratio=169]{beamer}

\usepackage{soul}
\usepackage{tikz}
\usepackage{multicol}
%\usepackage{gensymb}
\usepackage[utf8]{inputenc}

\usetikzlibrary{arrows.meta}

%\usepackage{hyperref}
\hypersetup{
    colorlinks=true,
    linkcolor=blue,
    filecolor=magenta,      
    urlcolor=cyan,
    pdftitle={(KTXFI1EBNF) Physics I. Practice},
}

% ----------------------------------------------------------------------------------------------------------------------------------------

\title{\includegraphics[width=45pt]{oepecsetlogo.png}\\(KTXFI1EBNF) Physics I. Practice}

\author{Dr.~Gábor FACSKÓ, PhD\\\tiny Assistant Professor / Senior Research Fellow\\\textit{facsko.gabor@uni-obuda.hu}}

\institute{\Tiny Óbuda University, Faculty of Electrical Engineering, 1084 Budapest, Tavaszmező u. 17.\\Wigner Research Centre for Physics, Department of Space Physics and Space Technology, 1121 Budapest, Konkoly-Thege Miklós út 29–33.\\ \url{https://wigner.hu/~facsko.gabor}}

\date{March~23,~2026}

\logo{\includegraphics[height=0.9 true cm]{kekcsik.png}}

% ----------------------------------------------------------------------------------------------------------------------------------------

\begin{document}

\frame{\titlepage}

% ----------------------------------------------------------------------------------------------------------------------------------------

\begin{frame}[fragile,squeeze,allowframebreaks]
\frametitle{Exercises}
\begin{itemize}
\setlength{\parskip}{0pt}
\setlength{\itemsep}{0pt}
\item \textbf{Thermodynamics I:} Temperature scales, thermal expansion, thermometers. State variables, ideal gas, state changes, ideal gas equation of state, Boyle’s law, Gay-Lussac’s laws, universal (combined) gas law, heat, specific heat, molar heat, heat capacity.
\item \textbf{Thermodynamics II:} Expansion work, internal energy, the first law of thermodynamics, forms of the first law in different state changes.
\item \textbf{Thermodynamics III:} Cyclic processes, Carnot cycle, Carnot’s theorem, reduced heat, entropy, the second law of thermodynamics.
  
\pagebreak
\item Problems:
\begin{enumerate}
\setcounter{enumi}{6}   
\item With one mole of monoatomic, ideal gas at $27\,{}^{\circ}C$, we supply 8.31 Wh of heat while the pressure remains constant (slow process).
\begin{enumerate}
\item[(a)] What will be the achieved temperature?
\item[(b)] How much does the internal energy of the gas change?
\item[(c)] What is the volume work done by the gas?
\end{enumerate}

\item We isothermally compress oxygen, which has a volume of 500 liters, a pressure of $10^{6}\,Pa$, and a temperature of $30{}^{\circ}C$, to one-fifth of its original volume.
\begin{enumerate}
\item[(a)] How much work do we perform during compression?
\item[(b)] What is the amount of heat absorbed by the gas
\item[(c)] and the change in internal energy?
\item[(d)] What will be the pressure of the gas at the end of the process?
\end{enumerate}

\pagebreak
\begin{columns}
\begin{column}{0.5\linewidth}
\begin{center}
\begin{tikzpicture}[scale=0.8]
    % Tengelyek rajzolása határozott nyilakkal a végükön
    \draw[thick, -{Stealth[length=3mm]}] (0,0) -- (5,0) node[right] {$V \, [m^3]$};
    \draw[thick, -{Stealth[length=3mm]}] (0,0) -- (0,5) node[above] {$p \, [Pa]$};

    % Pontok koordinátái
    \coordinate (state1) at (1, 1.5);
    \coordinate (state2) at (4, 4);

    % Az állapotokat összekötő barna nyíl
    \draw[brown, line width=1.2pt, -{Stealth[length=3mm]}] (state1) -- (state2);

    % Piros karikák (állapotok)
    \draw[red, fill=white, thick] (state1) circle (3pt) node[below left, black] {1};
    \draw[red, fill=white, thick] (state2) circle (3pt) node[above right, black] {2};

\end{tikzpicture}
\end{center}
\end{column}
\begin{column}{0.5\linewidth}
\item We transport a certain amount of air from state 1 to state 2 according to the figure. Does the temperature of the air increase or decrease during the process? How much work does the gas perform during the process if $p_{1} = 10^{5}\,Pa$, $p_{2} = 1.4 \cdot 10^{5}\,Pa$, $V_{1} = 0.07\,m^{3}$, $V_{2} = 0.11\,m^{3}$?
\end{column}
\end{columns}

\end{enumerate}

\end{itemize}

\end{frame}

% ----------------------------------------------------------------------------------------------------------------------------------------

\begin{frame}
%\frametitle{Minden jó, ha a vége jó}

\vfill
\begin{center}
\textbf{\huge The End}

\bigskip
Thank you for your attention!
\end{center}
\vfill
%\textbf{Acknowledgements} Thank you for the course materials for Szilárd~ZSÓKA and Zoltán~VARGA.
\end{frame}

\end{document}
